Formalization is in progress. The companion uses Mathlib's actual upper-half-plane action, complex differentiation, and real bilinear forms. The shared derivative theorem proves a \texttt{HasDerivAt} statement for the fractional transformation, not merely a formula named a derivative. Together with the previously proved action formula, this supplies the derivative used in the geometric calculations.

For \hyperref[hyperbolic.invariant-geometry]{uniqueness of the metric}, the formal field $B$ consists of real bilinear forms on $\mathbb C$, with explicit symmetry, positive-definiteness, and invariance under the actual fractional derivative. Its coefficients are proved to be $\lambda\operatorname{Re}(v\overline w)/y^2$ with $\lambda>0$. The proof follows the prose: the derivative of $k(\pi/4)$ is multiplication by $i$, which forces equal diagonal coefficients and zero mixed coefficient at $i$; the derivative of $s(x+iy)$ then transports this form to height $y$. Smoothness is unnecessary for this uniqueness direction and is not assumed in the Lean theorem. The candidate is now constructed as an actual real bilinear form, with symmetry, positive-definiteness and invariance proved; these facts supply existence at the level of coordinate forms. The smooth tangent-bundle interface remains open.

The normalized squared tangent length is $\operatorname{normSq}(v)/y^2$. Lean proves positivity for nonzero vectors, Euclidean normalization at $i$, and invariance under the derivative. Lean also proves invariance of the unsquared tangent length and of the full tangent inner product, so the substitution of a path velocity uses the same pointwise identity as the expanded prose. These are tangent statements; the chain rule along a piecewise smooth path, its length integral, and the infimum of path lengths remain unformalized.

For \hyperref[hyperbolic.invariant-density]{invariant measures with smooth densities}, Lean proves covariance of the positive density $1/y^2$ under the squared complex derivative and uniqueness up to the value at $i$ for any positive function satisfying that covariance law. The hypothesis is the pointwise law $\rho(gz)|g'(z)|^2=\rho(z)$. This is not yet a theorem about invariant measures: the real-Jacobian formula, change of variables, the passage from almost-everywhere equality of smooth densities to pointwise equality, and identification with induced Riemannian area are remaining obligations. The general Radon-measure uniqueness proof is in Section 1.8 and has not been formalized.

The manuscript now separates \hyperref[hyperbolic.vertical-geodesics]{vertical geodesics}, the \hyperref[hyperbolic.geodesics]{geodesic classification}, and the \hyperref[hyperbolic.distance-formula]{distance formula}. For the distance formula, Lean proves invariance of $J(z,w)=|z-w|^2/(\operatorname{Im}z\operatorname{Im}w)$ by the same cancellation of denominators as the prose. The circle-center construction, its straightening map, minimizing arcs, and the path-distance formula remain open. The paired SVG/PDF diagram is generated from one explicit circle and its M\"obius image; it is explanatory artwork, not proof evidence.

The \hyperref[hyperbolic.vertical-flow]{unit-speed example} now spells out the derivative of $ie^t$, the cancellation of Euclidean speed against height, and transport to $h^{-1}(ie^t)$ on a semicircular geodesic. The checked tangent-length invariance is the pointwise ingredient for this transport. The exponential path's derivative, its unit-speed specialization, the path chain rule, and the arclength/distance conclusions are not yet companion theorems; this clarification does not extend the recorded Lean coverage.

The gradient coefficient formula is checked through duality with an arbitrary real linear differential on the tangent coordinates: the constructed vector pairs with every tangent vector to give that differential. Identifying it with the differential of a smooth function and developing its smooth vector field remain further steps. Divergence is now defined as the unique continuous function satisfying integration by parts against every real smooth compactly supported test function, using the measure $\mu$. The prose proves existence by coordinate integration by parts and uniqueness by a test function supported where a hypothetical difference has one sign. This replaces the local-deformation derivation without changing the coordinate formula. Formalizing this characterization, its existence and uniqueness, and its compatibility with the gradient remains open. The Laplace--Beltrami formula, commutation with every smooth Riemannian isometry (including orientation reversal), and the Green identity also remain prose proofs with explicit formalization obligations; no Lean theorem about these analytic operators is claimed. Self-adjoint domains are deferred to the spectral chapter.

The \hyperref[hyperbolic.laplacian-isometries]{isometry-commutation proof} now applies the divergence characterization directly to the Laplacian. It first proves area preservation by cancellation of Euclidean area scaling with the hyperbolic density, then proves the gradient transformation rule. Integration by parts and change of variables identify the two candidate Laplacians against all compactly supported test functions; continuity gives pointwise equality. This retains every smooth Riemannian isometry, including orientation reversal, without introducing a transported vector field. These analytic steps remain formalization obligations; the existing tangent-form and gradient-duality theorems provide only their pointwise algebraic ingredients.

Further targets are hyperbolic path length and its infimum, existence and uniqueness of minimizing arcs, the exact distance identity, and those differential and integral operator statements. A formal distance formula alone would not prove the path-length characterization. Tangent products here are real; the later complex $L^2(\mathbb H,\mu)$ product is linear in its first argument.

Machine evidence for the implemented portions is recorded in \texttt{lean/audits/chapter01-release-check-2026-11-12/record.json}, against its source hashes. The section remains in progress, and Steven's correspondence review is pending independently of compilation and the axiom audit.

The horizontal-segment length formula \eqref{eq:horizontal-segment-length} now displays the constant-speed integral and is cited directly by the cusp discussion. It is the same geometric calculation; the current companion does not yet formalize this path integral.
